\documentclass[12pt,a4paper,oneside]{article}
\usepackage[utf8]{inputenc}
\usepackage[spanish,es-noshorthands]{babel}
\usepackage[T1]{fontenc}
\usepackage{lmodern}
\usepackage[margin=2.5cm]{geometry}
\usepackage{amsmath,amssymb}
\usepackage{xcolor}
\usepackage{hyperref}
\usepackage{graphicx}
\usepackage{booktabs}
\usepackage{tabularx}
\usepackage{enumitem}
\usepackage{multirow}
\usepackage{colortbl}
\usepackage{fancyhdr}
\usepackage{lastpage}
\usepackage{longtable}
\usepackage{listings}

\hypersetup{colorlinks=true,urlcolor=blue,linkcolor=black}

\definecolor{darkblue}{RGB}{31,56,100}
\definecolor{accent}{RGB}{79,143,247}
\definecolor{success}{RGB}{16,185,129}
\definecolor{warning}{RGB}{245,158,11}
\definecolor{danger}{RGB}{239,68,68}
\definecolor{lightgray}{RGB}{240,242,247}
\definecolor{codebg}{RGB}{248,249,252}
\definecolor{codekw}{RGB}{109,40,217}
\definecolor{codestr}{RGB}{16,120,40}
\definecolor{codecom}{RGB}{120,120,120}

\lstset{
  basicstyle=\ttfamily\small,
  backgroundcolor=\color{codebg},
  keywordstyle=\color{codekw}\bfseries,
  stringstyle=\color{codestr},
  commentstyle=\color{codecom}\itshape,
  showstringspaces=false,
  breaklines=true,
  frame=single,
  rulecolor=\color{lightgray},
  numbers=left,
  numberstyle=\tiny\color{codecom},
  xleftmargin=2em,
  framexleftmargin=1.5em
}

\pagestyle{fancy}
\fancyhf{}
\rhead{\textcolor{accent}{ISO SECURE}}
\lhead{\textcolor{accent}{Etapa 4 - Pruebas y Despliegue}}
\cfoot{Página \thepage\ de \pageref{LastPage}}
\renewcommand{\headrulewidth}{0.5pt}
\renewcommand{\headrule}{\color{accent}\hrule}

\title{
  {\Huge \textbf{ISO SECURE}}\\
  {\LARGE Plataforma de Cumplimiento ISO 27001}\\[2ex]
  {\Large \textcolor{accent}{Etapa 4: Pruebas y Despliegue}}\\[1ex]
  {\small Actividad Formativa 4 - 5\%}
}
\author{Agustín Peralta --- Universidad de San Buenaventura}
\date{Entrega: 18 de Mayo de 2026}

\begin{document}

\maketitle

\begin{abstract}
Este documento presenta las pruebas de seguridad dinámicas (DAST) y la configuración segura aplicadas en el proyecto ISO SECURE durante la fase de pruebas y despliegue. Cubre la estrategia DAST con OWASP ZAP, Nuclei y pruebas de fuzzing; el hardening de infraestructura (Nginx, Docker, PostgreSQL, sistema operativo); las cabeceras de seguridad HTTP; la gestión de secretos en producción; la estrategia de backups y los procedimientos de respuesta a incidentes. Las prácticas se alinean con OWASP ASVS 4.0, CIS Benchmarks y los controles A.8 (controles tecnológicos) y A.5 (organizacionales) de ISO/IEC 27001:2022.
\end{abstract}

\tableofcontents
\newpage

% ============================================================================
\section{Introducción}

La Etapa 4 cierra el ciclo de construcción con dos focos complementarios:

\begin{enumerate}[leftmargin=*]
  \item \textbf{Análisis Dinámico de Seguridad (DAST):} ejecutar la aplicación y atacarla deliberadamente para descubrir vulnerabilidades que el análisis estático no puede detectar (problemas de configuración, lógica de negocio, sesiones, autenticación).
  \item \textbf{Configuración Segura:} endurecer cada capa del stack (sistema operativo, contenedores, servidor web, base de datos, aplicación) siguiendo CIS Benchmarks y guías oficiales del proveedor.
\end{enumerate}

A diferencia del SAST (Etapa 3), que analiza código fuente, el DAST analiza el sistema corriendo. Ambos son complementarios y obligatorios para una cobertura de seguridad madura (DevSecOps).

% ============================================================================
\section{Análisis Dinámico (DAST)}

\subsection{Herramientas seleccionadas}

\begin{table}[h]
\centering
\small
\begin{tabularx}{\textwidth}{lXl}
\toprule
\textbf{Herramienta} & \textbf{Propósito} & \textbf{Tipo} \\
\midrule
OWASP ZAP & Scan automatizado + proxy interceptor + spider & DAST general \\
Nuclei & Plantillas YAML para CVEs conocidas y misconfigs & DAST plantillas \\
sqlmap & Pruebas dirigidas de inyección SQL & Especializado \\
nikto & Server scan: archivos comunes, headers, versiones & Web server \\
testssl.sh & Auditoría de TLS (ciphers, versiones, certs) & TLS/SSL \\
ffuf & Fuzzing de endpoints y parámetros & Fuzzing \\
JWT\_Tool & Pruebas específicas contra tokens JWT & API/Auth \\
Burp Suite Community & Análisis manual asistido & Manual \\
\bottomrule
\end{tabularx}
\end{table}

\subsection{Estrategia DAST por capas}

\subsubsection{Capa 1 — Reconocimiento (sin login)}

\begin{itemize}[leftmargin=*]
  \item Spider de ZAP sobre la URL pública: descubrir endpoints expuestos.
  \item nikto: archivos sensibles (\texttt{.git}, \texttt{.env}, \texttt{backup.sql}).
  \item testssl.sh: validar TLS 1.3, sin cifrados débiles, HSTS.
\end{itemize}

\subsubsection{Capa 2 — Autenticación}

\begin{itemize}[leftmargin=*]
  \item Enumeración de usuarios vía login (mensajes diferenciados).
  \item Brute force con throttle/lockout.
  \item Resistencia a credential stuffing (Have I Been Pwned).
  \item Validación de JWT: firma, alg=none, kid traversal, exp, nbf.
\end{itemize}

\subsubsection{Capa 3 — Autorizado (con cuenta de prueba)}

\begin{itemize}[leftmargin=*]
  \item ZAP Active Scan con contexto autenticado (script de re-auth).
  \item Tests de IDOR: cambiar IDs en URL/body.
  \item Tests multi-tenant: con token de empresa A intentar leer/escribir en empresa B.
  \item Privilege escalation: con rol supervisor intentar acciones de admin.
\end{itemize}

\subsubsection{Capa 4 — Fuzzing}

\begin{itemize}[leftmargin=*]
  \item ffuf sobre parámetros JSON con payloads de OWASP juice-shop.
  \item Pruebas de límites: strings de 1MB, enteros negativos, fechas inválidas.
  \item Inyecciones: SQL, NoSQL, comandos, plantilla (SSTI), LDAP, XPath.
\end{itemize}

\subsection{Plan de ejecución}

\begin{lstlisting}[language=bash]
# 1. Scan baseline ZAP en CI (cada PR)
docker run --rm -v $(pwd):/zap/wrk owasp/zap2docker-stable \
  zap-baseline.py -t https://staging.iso-secure.app -r zap-report.html

# 2. Nuclei contra staging diariamente
nuclei -u https://staging.iso-secure.app \
  -t cves/ -t misconfiguration/ -t exposures/ \
  -severity critical,high -o nuclei-report.txt

# 3. testssl.sh contra produccion (semanal)
testssl.sh --severity HIGH https://iso-secure.app

# 4. ZAP autenticado (release)
zap-full-scan.py -t https://staging.iso-secure.app \
  -z "-config replacer.full_list(0).description=auth \
       -config replacer.full_list(0).enabled=true \
       -config replacer.full_list(0).matchtype=REQ_HEADER \
       -config replacer.full_list(0).matchstr=Authorization \
       -config replacer.full_list(0).replacement=Bearer $TOKEN"
\end{lstlisting}

\subsection{Catálogo esperado de hallazgos y mitigaciones}

\begin{longtable}{p{1cm}p{4cm}cp{6cm}}
\toprule
\textbf{ID} & \textbf{Hallazgo} & \textbf{Sev.} & \textbf{Mitigación} \\
\midrule
\endhead
D-01 & Cookies sin flag HttpOnly/Secure & Alta & Configurar atributos cookie + SameSite=Lax \\
D-02 & CORS demasiado permisivo (\texttt{*}) & Alta & Whitelist explícita por entorno \\
D-03 & Headers de seguridad faltantes (CSP, X-Frame) & Media & Añadir en Nginx (sección Hardening) \\
D-04 & Mensajes de error con stack trace en prod & Media & \texttt{DEBUG=False} + handler genérico \\
D-05 & Robots.txt expone rutas internas & Baja & Eliminar referencias internas \\
D-06 & Server: nginx/1.x revela versión & Baja & \texttt{server\_tokens off;} \\
D-07 & TLS 1.0/1.1 habilitado & Alta & Solo TLS 1.2+ (ideal 1.3) \\
D-08 & Cifrados RC4/3DES soportados & Alta & Suite de cifrados modernos (Mozilla intermediate) \\
D-09 & Brute force sin lockout & Alta & Rate limit + lockout temporal en Supabase \\
D-10 & JWT alg=none aceptado & Crítica & Validar contra Supabase /auth/v1/user (ya implementado) \\
D-11 & IDOR en \texttt{/incidents/\{id\}} & Crítica & Verificar ownership en service layer \\
D-12 & Cross-tenant via parámetro & Crítica & Helper \texttt{\_empresa\_filter} (ya implementado) \\
D-13 & SSRF en webhooks n8n & Alta & Allowlist de hosts salientes \\
D-14 & Open redirect en /login?return\_to & Media & Validar dominio destino \\
D-15 & Falta de rate limit global & Alta & slowapi + Nginx limit\_req \\
\bottomrule
\end{longtable}

% ============================================================================
\section{Configuración Segura — Hardening por capa}

\subsection{Sistema Operativo (host)}

\begin{itemize}[leftmargin=*]
  \item Distribución mínima (Debian slim o Alpine).
  \item Actualizaciones automáticas de seguridad (\texttt{unattended-upgrades}).
  \item Usuario no-root para correr servicios.
  \item SSH: solo claves, sin password, puerto no estándar, fail2ban.
  \item Firewall ufw: permitir 22, 80, 443; denegar resto.
  \item Auditd habilitado para eventos sensibles.
  \item Sin servicios innecesarios; \texttt{systemd-analyze} review.
\end{itemize}

\subsection{Docker / Contenedores}

\begin{itemize}[leftmargin=*]
  \item Imágenes base oficiales y versionadas (no \texttt{latest}).
  \item Multi-stage builds para reducir superficie de ataque.
  \item Usuario no-root en el \texttt{Dockerfile} (\texttt{USER 10001}).
  \item Capabilities mínimas (\texttt{--cap-drop=ALL}).
  \item Read-only filesystem cuando sea posible.
  \item Healthchecks declarados.
  \item Escaneo con Trivy antes de promover a producción.
\end{itemize}

\textbf{Ejemplo Dockerfile hardened:}

\begin{lstlisting}[language=bash]
FROM python:3.13-slim AS builder
WORKDIR /app
COPY requirements.txt .
RUN pip install --user --no-cache-dir -r requirements.txt

FROM python:3.13-slim
RUN groupadd -r app && useradd -r -g app -u 10001 app
COPY --from=builder /root/.local /home/app/.local
COPY --chown=app:app . /app
WORKDIR /app
USER app
ENV PATH=/home/app/.local/bin:$PATH
HEALTHCHECK --interval=30s --timeout=3s CMD python -c "import urllib.request; urllib.request.urlopen('http://localhost:8000/health')"
CMD ["uvicorn", "main:app", "--host", "0.0.0.0", "--port", "8000"]
\end{lstlisting}

\subsection{Nginx (reverse proxy + TLS termination)}

\begin{itemize}[leftmargin=*]
  \item TLS 1.2+ obligatorio; 1.3 preferido.
  \item Suite de cifrados Mozilla intermediate.
  \item HSTS con \texttt{max-age=31536000; includeSubDomains; preload}.
  \item \texttt{server\_tokens off;} para ocultar versión.
  \item Rate limiting por IP en endpoints sensibles (login).
  \item Buffer sizes para mitigar Slowloris.
  \item Logs en formato JSON con IP, user-agent, status, latencia.
\end{itemize}

\textbf{Headers de seguridad obligatorios:}

\begin{lstlisting}[language=bash]
# nginx headers (extracto)
add_header Strict-Transport-Security "max-age=31536000; includeSubDomains; preload" always;
add_header X-Content-Type-Options "nosniff" always;
add_header X-Frame-Options "DENY" always;
add_header Referrer-Policy "strict-origin-when-cross-origin" always;
add_header Permissions-Policy "geolocation=(), microphone=(), camera=()" always;
add_header Content-Security-Policy "default-src 'self'; script-src 'self' 'unsafe-inline'; style-src 'self' 'unsafe-inline'; img-src 'self' data: https:; connect-src 'self' https://*.supabase.co; frame-ancestors 'none';" always;
\end{lstlisting}

\subsection{PostgreSQL (Supabase managed)}

Aunque Supabase gestiona muchos controles, las prácticas que aplicamos del lado cliente son:

\begin{itemize}[leftmargin=*]
  \item Solo conexiones SSL (\texttt{sslmode=require}).
  \item Roles separados: \texttt{app\_user} (CRUD) y \texttt{migrator} (DDL).
  \item Sin \texttt{SUPERUSER} para conexiones de aplicación.
  \item Row Level Security (RLS) habilitado en tablas multi-tenant.
  \item Backups automáticos diarios + Point-in-Time Recovery 7 días.
  \item Encriptación at-rest activa por defecto.
  \item Audit logging habilitado para DDL y privilegios.
\end{itemize}

\subsection{Aplicación FastAPI}

\begin{itemize}[leftmargin=*]
  \item \texttt{DEBUG=False} y \texttt{docs\_url=None} en producción (o tras auth admin).
  \item Workers Uvicorn $\geq 2 \times$ vCPU para tolerancia.
  \item Timeouts: \texttt{--timeout-keep-alive=5}, \texttt{--limit-concurrency=1000}.
  \item Middleware de logging estructurado con request\_id.
  \item Trusted Hosts middleware con whitelist explícita.
  \item GZip middleware con \texttt{minimum\_size} para evitar BREACH.
\end{itemize}

\subsection{Frontend React (build estático)}

\begin{itemize}[leftmargin=*]
  \item Source maps solo en staging; deshabilitados en prod.
  \item Subresource Integrity (SRI) en scripts externos.
  \item Variables sólo \texttt{VITE\_} (públicas por diseño); nada sensible en bundle.
  \item Sentry/Posthog desactivados en preview branches.
\end{itemize}

% ============================================================================
\section{Gestión de Secretos en Producción}

\subsection{Principios}

\begin{enumerate}[leftmargin=*]
  \item \textbf{Nunca en el código.} Detección con \texttt{detect-secrets} en pre-commit y CI.
  \item \textbf{Nunca en logs.} Filtros de redacción en el logger.
  \item \textbf{Rotación periódica.} 90 días para credenciales operativas.
  \item \textbf{Mínimo privilegio.} Cada servicio con sus propias credenciales.
  \item \textbf{Cifrado en reposo.} Vault del proveedor (Supabase Vault, GitHub Secrets, Coolify Secrets).
\end{enumerate}

\subsection{Inventario de secretos}

\begin{table}[h]
\centering
\small
\begin{tabularx}{\textwidth}{lXl}
\toprule
\textbf{Secreto} & \textbf{Uso} & \textbf{Rotación} \\
\midrule
DATABASE\_URL & Conexión PostgreSQL & 90 días \\
SUPABASE\_URL & Endpoint del proveedor & estática \\
SUPABASE\_ANON\_KEY & Auth cliente (pública por diseño) & 180 días \\
SUPABASE\_SERVICE\_ROLE\_KEY & Operaciones server-side & 90 días \\
JWT\_VALIDATION\_URL & Endpoint de validación & estática \\
N8N\_WEBHOOK\_SECRET & HMAC de webhooks salientes & 90 días \\
SMTP\_PASSWORD & Emails transaccionales & 90 días \\
\bottomrule
\end{tabularx}
\end{table}

% ============================================================================
\section{Estrategia de Despliegue}

\subsection{Pipeline de promoción}

\begin{center}
\begin{tabular}{|c|c|c|c|}
\hline
\rowcolor{accent!20}
\textbf{Entorno} & \textbf{Trigger} & \textbf{Pruebas} & \textbf{Aprobación} \\
\hline
PR preview & abrir PR & smoke + ZAP baseline & automática \\
\hline
staging & merge a \texttt{develop} & full ZAP + e2e Playwright & automática \\
\hline
producción & merge a \texttt{main} & smoke post-deploy & manual (1 reviewer) \\
\hline
\end{tabular}
\end{center}

\subsection{Estrategia blue-green simplificada (Coolify)}

\begin{itemize}[leftmargin=*]
  \item Coolify levanta nueva versión del contenedor.
  \item Healthcheck \texttt{/health} debe pasar.
  \item Si pasa, swap de tráfico en el reverse proxy interno.
  \item Si falla, rollback automático a la versión previa (último commit verde).
\end{itemize}

\subsection{Plan de rollback}

\begin{enumerate}[leftmargin=*]
  \item Detectar regresión vía healthcheck, monitoreo o reporte de usuario.
  \item Coolify: re-deploy del último commit etiquetado como \texttt{stable}.
  \item Si involucra migración de DB: ejecutar \texttt{alembic downgrade -1}.
  \item Post-mortem y RCA dentro de 48 hrs (control A.5.27).
\end{enumerate}

% ============================================================================
\section{Backups y Continuidad}

\begin{table}[h]
\centering
\small
\begin{tabularx}{\textwidth}{lXl}
\toprule
\textbf{Activo} & \textbf{Estrategia} & \textbf{RPO / RTO} \\
\midrule
Base de datos & PITR Supabase 7 días + snapshot diario externo & 1h / 2h \\
Código & Git remoto (GitHub) + mirror en repo URPE & 0 / 30min \\
Configuración (Coolify) & Export semanal del project a S3 cifrado & 24h / 1h \\
Secretos & Backup cifrado del vault en cold storage & 24h / 4h \\
Logs & Retención 90d + agregación off-site & 1h / N/A \\
\bottomrule
\end{tabularx}
\end{table}

% ============================================================================
\section{Respuesta a Incidentes}

\subsection{Niveles de severidad}

\begin{table}[h]
\centering
\small
\begin{tabularx}{\textwidth}{lXl}
\toprule
\textbf{Nivel} & \textbf{Descripción} & \textbf{SLA respuesta} \\
\midrule
P0 (Crítico) & Brecha activa, indisponibilidad total, fuga de datos & 15 min \\
P1 (Alto) & Vulnerabilidad explotable confirmada, degradación severa & 1 h \\
P2 (Medio) & Vulnerabilidad sin exploit público, función no crítica caída & 4 h \\
P3 (Bajo) & Hardening pendiente, hallazgo informativo & 5 días hábiles \\
\bottomrule
\end{tabularx}
\end{table}

\subsection{Procedimiento IR (control A.5.24-A.5.28)}

\begin{enumerate}[leftmargin=*]
  \item \textbf{Detección} → alerta automática o reporte manual.
  \item \textbf{Triaje} → asignar severidad y propietario.
  \item \textbf{Contención} → aislar (revocar tokens, bloquear IP, deshabilitar usuario, blue-green rollback).
  \item \textbf{Erradicación} → corregir causa raíz, parchear, actualizar.
  \item \textbf{Recuperación} → restaurar servicio + validar integridad.
  \item \textbf{Lecciones aprendidas} → post-mortem + actualización de catálogo de amenazas (Etapa 2).
\end{enumerate}

% ============================================================================
\section{Monitoreo y Observabilidad}

\begin{itemize}[leftmargin=*]
  \item \textbf{Logs:} salida JSON estructurada → agregador (Loki/Grafana o Better Stack).
  \item \textbf{Métricas:} CPU, RAM, latencia p95, error rate (Prometheus + Grafana).
  \item \textbf{Traces:} OpenTelemetry para request flow end-to-end.
  \item \textbf{Alertas:} error rate $> 1\%$ por 5min, latencia p95 $> 2$s, healthcheck fallido.
  \item \textbf{Dashboard ejecutivo:} disponibilidad mensual (SLA 99.5\%), MTTR, número de incidentes por severidad.
\end{itemize}

% ============================================================================
\section{Mapeo a ISO 27001:2022}

\begin{table}[h]
\centering
\small
\begin{tabularx}{\textwidth}{lX}
\toprule
\textbf{Control Anexo A} & \textbf{Implementación} \\
\midrule
A.5.24 Planificación de IR & Procedimiento documentado + SLAs por severidad \\
A.5.25 Evaluación de eventos & Triaje formal en cada alerta \\
A.5.27 Lecciones aprendidas & Post-mortem en 48h + retroalimentación al modelo de amenazas \\
A.5.30 Continuidad TIC & RPO/RTO definidos + PITR + mirror de código \\
A.8.7 Protección contra malware & Trivy en CI + actualizaciones de base OS \\
A.8.8 Gestión de vulnerabilidades técnicas & DAST + Nuclei + Dependabot \\
A.8.9 Gestión de la configuración & Coolify + IaC + secretos en vault \\
A.8.13 Backup de información & Backups diarios + PITR \\
A.8.15 Logging & Logs JSON estructurados + retención 90d \\
A.8.16 Monitoreo de actividades & Prometheus + alertas \\
A.8.20 Seguridad de redes & TLS 1.3 + firewall + rate limit \\
A.8.23 Filtrado web & CSP estricto + CORS allowlist \\
A.8.26 Requisitos de seguridad de aplicaciones & DAST regular como gate de release \\
A.8.27 Arquitectura segura & Hardening por capa documentado \\
A.8.28 Codificación segura & Resultados DAST retroalimentan estándares de código \\
\bottomrule
\end{tabularx}
\end{table}

% ============================================================================
\section{Conclusiones}

La Etapa 4 cierra el ciclo de aseguramiento técnico previo a la entrega final:

\begin{enumerate}[leftmargin=*]
  \item \textbf{DAST automatizado} en CI (ZAP baseline en cada PR; Nuclei diario; full ZAP autenticado por release).
  \item \textbf{Hardening por capa} documentado: SO, Docker, Nginx, PostgreSQL, FastAPI, frontend.
  \item \textbf{Gestión de secretos} con vault, rotación 90 días y detección automática en pre-commit/CI.
  \item \textbf{Despliegue blue-green} mediante Coolify con rollback automático.
  \item \textbf{Backups y continuidad} con RPO 1h / RTO 2h en la base de datos.
  \item \textbf{Respuesta a incidentes} con SLA por severidad y post-mortem obligatorio.
\end{enumerate}

Las pruebas dinámicas no son un evento puntual: son un proceso continuo que retroalimenta al modelo de amenazas (Etapa 2) y a las prácticas de codificación segura (Etapa 3), cumpliendo el ciclo PDCA exigido por ISO 27001:2022.

\section*{Referencias}
\begin{itemize}[leftmargin=*]
  \item OWASP Web Security Testing Guide (WSTG) v4.2.
  \item OWASP Application Security Verification Standard (ASVS) 4.0.3.
  \item CIS Docker Benchmark v1.6 y CIS Nginx Benchmark v2.0.
  \item Mozilla SSL Configuration Generator (Intermediate Profile).
  \item NIST SP 800-61 Rev. 2 — Computer Security Incident Handling Guide.
  \item ISO/IEC 27001:2022 — Anexo A.5 y A.8.
\end{itemize}

\end{document}
